\section{La diffusion linéaire : un modèle de référence}
\label{sec:heat}

\subsection{Évolution, flux et conditions aux limites}

Le modèle de chaleur part de l'observation bruitée $f$ et recherche une
famille d'images $u(\cdot,t)$ vérifiant
\begin{equation}
 \partial_t u=\Delta u \quad\text{dans }\Omega,\qquad
 u(\cdot,0)=f,\qquad \partial_{\boldsymbol n}u=0
 \quad\text{sur }\partial\Omega.
 \label{eq:heat-continuous}
\end{equation}
Le Laplacien $\Delta u=\partial_{xx}u+\partial_{yy}u$ mesure la différence
entre une valeur et ses environs immédiats : une pointe isolée tend à
diminuer et un creux isolé tend à augmenter. L'équation ne distingue pas
l'origine de ces variations. Elle atténue donc une fluctuation de bruit,
mais également une transition appartenant réellement à l'image.
Le coefficient de diffusion est fixé à un ; toute modification de ce
coefficient changerait l'échelle temporelle de comparaison.

En introduisant le flux physique $\boldsymbol J=-\nabla u$, on écrit
$\partial_t u+\operatorname{div}\boldsymbol J=0$.
La condition $\partial_{\boldsymbol n}u=0$ impose
$\boldsymbol J\cdot\boldsymbol n=0$ : aucune quantité n'est échangée avec
l'extérieur. Pour une solution suffisamment régulière, l'intégration
sur le rectangle donne
\begin{equation}
 \frac{\dd}{\dd t}\int_\Omega u\,\dd x\,\dd y
 =\int_{\partial\Omega}\nabla u\cdot\boldsymbol n\,\dd S=0.
 \label{eq:heat-mass}
\end{equation}
Ainsi, la moyenne spatiale reste celle de l'observation. En particulier,
le débruitage ne supprime pas la moyenne du bruit effectivement tiré :
même si son espérance statistique est nulle, sa moyenne sur un échantillon
fini n'est généralement pas exactement nulle.

\subsection{Échelles et fréquences}

Sur l'espace entier $\R^2$, pour une donnée initiale dans un cadre où la
convolution est bien définie, par exemple $f\in L^1(\R^2)\cap L^2(\R^2)$,
la solution à $t>0$ s'exprime par
\begin{equation}
 u(x,y,t)=(G_t*f)(x,y),\qquad
 G_t(z)=\frac{1}{4\pi t}\exp\left(-\frac{|z|^2}{4t}\right).
 \label{eq:heat-gaussian}
\end{equation}
La variance de ce noyau est $2t$ dans chacune des deux directions.
Avec la convention angulaire pour la transformation de Fourier,
on obtient le multiplicateur
$\exp(-t|\xi|^2)$ : les composantes de grande fréquence spatiale sont
atténuées plus rapidement. Cette représentation éclaire l'intérêt de la
chaleur pour un bruit fluctuant à l'échelle du pixel. Elle explique aussi
pourquoi une frontière nette s'élargit avec le temps.

Le problème effectivement étudié n'est toutefois pas posé sur $\R^2$.
Sur un rectangle avec conditions de Neumann, la solution est engendrée
par le Laplacien de Neumann ; sa représentation utilise ses modes propres
ou un noyau adapté aux frontières. Des réflexions de la donnée peuvent
permettre une représentation équivalente dans un cadre approprié, mais
la convolution de la seule image finie avec $G_t$ sur tout l'espace
n'est pas, sans traitement du bord, la solution exacte de
\eqref{eq:heat-continuous}. Une extension par zéro imposerait notamment
un comportement différent aux frontières. Cette distinction est
essentielle pour relier le modèle au code conservatif.

\subsection{Deux dissipations et leurs hypothèses}

Pour une solution régulière, l'énergie de Dirichlet
\begin{equation}
 E_D(u)=\frac12\int_\Omega |\nabla u|^2\,\dd x\,\dd y
 \label{eq:dirichlet-energy}
\end{equation}
mesure la variation spatiale. Si une perturbation $v$ respecte le cadre
variationnel de Neumann, son premier accroissement s'écrit
\begin{equation}
 \dd E_D(u)[v]
 =\int_\Omega \nabla u\cdot\nabla v
 =-\int_\Omega (\Delta u)v,
 \label{eq:dirichlet-firstvariation}
\end{equation}
le terme de bord étant nul. L'équation de chaleur est donc formellement
le flot de gradient négatif de $E_D$ pour la métrique $L^2$.
Pour une solution assez régulière pour dériver l'énergie et effectuer
l'intégration par parties, ce calcul donne
\begin{equation}
 \frac{\dd}{\dd t}E_D(u(t))
 =-\int_\Omega (\Delta u)^2\,\dd x\,\dd y\leq0.
 \label{eq:dirichlet-dissipation}
\end{equation}
Ce sont des identités pour le modèle continu régulier ; une donnée
pixelisée et bruitée n'est pas supposée posséder initialement toutes
les dérivées utilisées. Elles motivent le modèle, sans remplacer une
analyse du schéma entièrement discret.

Une seconde quantité, distincte de $E_D$, est la variance spatiale
quadratique. Avec $\bar f=|\Omega|^{-1}\int_\Omega f$, la conservation de
la moyenne et l'intégration par parties donnent
\begin{equation}
 \frac{\dd}{\dd t}\frac12\int_\Omega (u-\bar f)^2
 =-\int_\Omega |\nabla u|^2\leq0.
 \label{eq:heat-variance}
\end{equation}
La chaleur réduit les variations par rapport à la moyenne, et non
nécessairement l'erreur par rapport à $\utrue$ à chaque instant. Lorsque
les détails de l'image propre sont suffisamment atténués, l'erreur de
reconstruction peut augmenter alors même que l'énergie continue de
décroître. Le sur-lissage ne constitue donc pas une instabilité du
modèle : il révèle un choix de temps inadapté à l'objectif de débruitage.

Cette lecture en termes d'échelles et de diffusion constitue le point
de comparaison des modèles dépendant des variations locales
\cite{weickert1998}. Dans les expériences, le temps d'arrêt est par
conséquent un paramètre scientifique aussi important que le choix de
l'équation.
